\documentclass[aspectratio=169,professionalfonts]{beamer}
\usetheme{Madrid}
\usecolortheme{default}
\usepackage{amsmath,amssymb}
\usepackage[T1]{fontenc}
\title{STM-GARSOS Processing: Concise Overview}
\author{STM-GARSOS workflow}
\date{\today}
\begin{document}
\begin{frame}\titlepage\end{frame}
\begin{frame}{Objective and model}
\begin{itemize}
\item Represent each dynamic image as $X_t(x,y)=\sum_{\ell=1}^{L} {\rm ST}_{t,\ell}(x,y) C_\ell(x,y)$.
\item Estimate L spatial coefficient images from multicoil radial data.
\item Synthesize the complete dynamic image series after the coefficient solve.
\end{itemize}
\end{frame}
\begin{frame}{Raw-data preparation}
\begin{itemize}
\item Load GAR-SOS metadata, raw k-space, and a vendor-reconstructed reference volume.
\item Correct readout delay and prepare partition data.
\item Estimate/apply coil-noise whitening.
\item Generate normalized [kx, ky, kz] trajectory coordinates.
\item Compute radial density-compensation weights.
\end{itemize}
\end{frame}
\begin{frame}{Frame formation}
\begin{itemize}
\item Flattened arrays retain one row per readout/partition/spoke sample.
\item Spokes are grouped consecutively into temporal frames.
\item Historical baseline: 15 spokes per frame and approximately 133 complete frames.
\item Incomplete trailing spokes are discarded.
\end{itemize}
\end{frame}
\begin{frame}{Calibration-data path}
\begin{itemize}
\item Each frame is reconstructed with 3-D Kaiser-Bessel gridding.
\item Oversampled grid: approximately 308 x 308 x 132.
\item 3-D inverse FFT and centered resize produce a 224 x 224 x 96 volume.
\item A selected z slice is transformed to 2-D Cartesian k-space.
\item A central 32 x 32 calibration square is extracted.
\end{itemize}
\end{frame}
\begin{frame}{STM calibration}
\begin{itemize}
\item $k_{\rm Cal}$ is supplied to $\texttt{STM\_computation}$.
\item The calibration data span the temporal frames.
\item The STM algorithm estimates ST maps and eigenvalues.
\item Historical baseline: L = 7 and 224 x 224 spatial map dimensions.
\end{itemize}
\end{frame}
\begin{frame}{Radial STM data}
\begin{itemize}
\item The historical baseline used a direct centered 58-point partition IFFT.
\item One z index is selected and the coil dimension is preserved.
\item Output radial data: [readout x spokes/frame x coils x frames].
\item In-plane kx/ky coordinates are retained for the NUFFT.
\end{itemize}
\end{frame}
\begin{frame}{Multicoil radial operator}
\begin{itemize}
\item Forward: $d_c={\cal N}_{\rm radial}(S_c \odot X)$.
\item Adjoint: ${\cal F}^H(d)=\sum_c S_c^* \odot {\cal N}_{\rm radial}^H(d_c)$.
\item Sensitivity maps are forward receive-coil maps.
\item Coils remain separate in the data fidelity term.
\end{itemize}
\end{frame}
\begin{frame}{STM operator and PCG}
\begin{itemize}
\item STM synthesis maps coefficients C to dynamic images X.
\item The radial operator maps X to predicted multicoil radial data.
\item Solve $({\cal A}^H{\cal A})C={\cal A}^H y$ with PCG.
\item Optional Tikhonov solve adds lambda I.
\end{itemize}
\end{frame}
\begin{frame}{Output and historical caveats}
\begin{itemize}
\item Output: coefficient images plus synthesized [224 x 224 x frames] image series.
\item The initial baseline used the direct 58-point z/partition IFFT.
\item Later work addresses forward-map estimation, z-grid consistency, scaling, and conditioning.
\item The slides describe the original stable path, not all later corrections.
\end{itemize}
\end{frame}
\end{document}
